\begin{figure}[t]
\centering
\resizebox{\columnwidth}{!}{%
\begin{tikzpicture}
\begin{axis}[
  width=8.6cm, height=5.7cm,
  ybar, bar width=7pt,
  symbolic x coords={heat, qw-L, qw-A, qw-Lsym, qw-multi},
  xtick=data, x tick label style={font=\scriptsize, rotate=18, anchor=east},
  ymin=0.05, ymax=0.099,
  ylabel={error (lower is better)},
  legend cell align=left,
  legend style={font=\scriptsize, at={(0.02,1.02)}, anchor=south west,
                legend columns=2, fill=white, draw=gray!40},
  ymajorgrids, grid style={gray!18},
  tick label style={font=\footnotesize}, label style={font=\small},
  axis line style={gray!50},
  enlarge x limits=0.13,
]
\addplot[draw=steel!80!black, fill=steel!35] coordinates {
  (heat,0.069) (qw-L,0.071) (qw-A,0.066) (qw-Lsym,0.071) (qw-multi,0.069)};
\addlegendentry{MAE}
\addplot[draw=amber!80!black, fill=palegold] coordinates {
  (heat,0.085) (qw-L,0.077) (qw-A,0.075) (qw-Lsym,0.085) (qw-multi,0.082)};
\addlegendentry{far-node AURC}
% highlight the best variant
\node[star, star points=5, star point ratio=2.3, fill=amber, draw=amber!70!black,
      inner sep=1.3pt] at (axis cs:qw-A,0.0905) {};
\node[font=\scriptsize\bfseries, text=amber!75!black] at (axis cs:qw-A,0.087) {best};
\end{axis}
\end{tikzpicture}}
\caption{Exp~B2, quantum-walk variant probe ($264$ satellites, $5$ seeds). The base
matrix of the walk matters: the adjacency walk (qw-A, starred) is the best operator on
both metrics at matched parameters, edging the classical heat kernel on the far-node
AURC, while the Laplacian (qw-L) and normalized-Laplacian (qw-Lsym) walks tie heat and
the higher-capacity multi-time variant does not improve on qw-A.}
\label{fig:expB2}
\end{figure}
